\documentclass[10pt]{article}
\usepackage[margin=0.8in]{geometry}
\usepackage{amsmath,amssymb,amsthm,mathtools}
\usepackage{graphicx}
\usepackage{xcolor}
\usepackage{microtype}
\usepackage[hidelinks]{hyperref}
\setlength{\parindent}{0pt}
\setlength{\parskip}{5pt}
\newtheorem{theorem}{Theorem}
\newtheorem{lemma}[theorem]{Lemma}
\newcommand{\R}{\mathbb R}
\newcommand{\I}{\mathcal I}
\newcommand{\E}{\mathcal E}
\newcommand{\Om}{(\Omega)}
\newcommand{\Cin}{C^\infty}
\newcommand{\ran}{\operatorname{ran}}

\title{Property \(\Om\) for Compacta on Finitely Many Parallel Lines\\
\large A higher-dimensional partial result for Conjecture 2.3 of arXiv:1611.06808}
\author{Automated proof attempt\\
\texttt{agent\_lane\_03} \quad model: GPT5.6}
\date{11 August 2026}

\begin{document}
\maketitle

\begin{center}
\fcolorbox{blue!45!black}{blue!4}{\parbox{0.92\linewidth}{
\textbf{Status: partial result, likely valid, pending human verification.}
The source proves \(\Om\) for vanishing ideals on the real line and
conjectures it for arbitrary compact subsets of \(\R^d\).  Here the conjecture
is proved in every dimension for arbitrary compact sets supported on finitely
many parallel affine lines.}}
\end{center}

\section*{Source question}

The excerpt is Conjecture 2.3 on printed page 6 of the source paper.

\begin{center}
\includegraphics[width=0.96\linewidth]{question_crop.png}
\end{center}

For a compact \(K\subset\R^d\), let
\[
 \I_K=\{f\in\Cin(\R^d): f|_K=0\}.
\]
The source proves that every closed ideal of \(\Cin(\R)\) has property
\(\Om\) \cite[Theorem 2.1]{FJW}.  The following transfers that theorem to a
genuinely higher-dimensional, geometrically broad class.

\begin{theorem}[Parallel-line case of Conjecture 2.3]
\label{thm:main}
Let \(d\ge2\), let \(v\in\R^d\) be a unit vector, and let
\(b_1,\ldots,b_q\in v^\perp\) be distinct.  Put
\[
 M=\bigcup_{i=1}^q(b_i+\R v).
\]
For every compact set \(K\subset M\), with no regularity or density
assumption on its line slices, the vanishing ideal \(\I_K\) in
\(\Cin(\R^d)\) has property \(\Om\).
\end{theorem}

\section*{Proof intuition}

Restriction to the carrier \(M\) has a very explicit continuous linear right
inverse: interpolate only in the transverse variable.  This splits all smooth
functions into a function vanishing on every carrier line plus independently
chosen smooth functions along the lines.  Imposing zero values on
\(K\) then acts only on those one-dimensional components.  The ambient ideal
therefore becomes a topological direct sum of a complemented smooth-function
subspace and finitely many one-dimensional vanishing ideals.  Property
\(\Om\) is stable under exactly these operations.

\section*{Proof}

We first record the two standard permanence inputs, including their short
proofs so that the argument does not depend on an unstated functional-analytic
fact.

\begin{lemma}\label{lem:permanence}
The Fr\'echet space \(\Cin(\R^d)\) has \(\Om\).  A complemented subspace of a
Fr\'echet space with \(\Om\) has \(\Om\), and a finite direct sum of spaces
with \(\Om\) has \(\Om\).
\end{lemma}

\begin{proof}
Use the seminorms
\[
 p_n(f)=\max_{|\alpha|\le n,\,|x|\le n}|\partial^\alpha f(x)|.
\]
Fix \(n\), choose a smooth cutoff \(\psi\) equal to one on a neighborhood
of the closed ball of radius \(n\) and supported in the ball of radius
\(m=n+1\), and let \(\rho_\delta\) be a standard mollifier at scale
\(0<\delta<\delta_0\), with \(\delta_0\) small enough for the chosen
neighborhood.  If \(p_m(f)\le1\), put
\(y=\rho_\delta*(\psi f)\).  Taylor's formula and the convolution estimates
give, for every \(k\ge m\),
\[
 p_n(f-y)\le C_n\delta,
 \qquad p_k(y)\le C_{n,k}\delta^{-k}.
\]
Taking \(\delta\) proportional to \(\varepsilon\) for
\(0<\varepsilon<1\), and \(y=0\) for \(\varepsilon\ge1\), is precisely the
approximation formulation of \(\Om\) recalled in the source immediately
before its Theorem 2.1.

If \(P:X\to Y\) is a continuous projection, apply an \(\Om\)-decomposition
in \(X\) to an element of \(Y\) and then apply \(P\) to both summands;
continuity merely shifts the seminorm indices and constants.  This proves
the complemented-subspace assertion.  Applying the defining estimates in
each coordinate, taking the maximum of the finitely many indices and
constants, proves the finite-direct-sum assertion.
\end{proof}

Write every \(x\in\R^d\) uniquely as \(x=tv+y\), where
\(t=\langle x,v\rangle\) and \(y\in v^\perp\).  Choose
\(\chi_1,\ldots,\chi_q\in\Cin(v^\perp)\) such that
\[
 \chi_i(b_j)=\delta_{ij}.
\]
For example, disjoint transverse bump functions suffice.  Define
\[
 R:\Cin(\R^d)\longrightarrow\bigoplus_{i=1}^q\Cin(\R),
 \qquad Rf=\bigl(t\mapsto f(tv+b_i)\bigr)_{i=1}^q,
\]
and
\[
 E:\bigoplus_{i=1}^q\Cin(\R)\longrightarrow\Cin(\R^d),
 \qquad E(g_1,\ldots,g_q)(tv+y)=\sum_{i=1}^q\chi_i(y)g_i(t).
\]
Both maps are continuous for the usual smooth compact-open Fr\'echet
topologies, and \(RE=I\).  Hence \(P=ER\) is a continuous projection,
\(\ker P=\ker R=\I_M\), and
\begin{equation}
 \Cin(\R^d)=\I_M\oplus E\!\left(\bigoplus_{i=1}^q\Cin(\R)\right)
 \label{eq:ambient-split}
\end{equation}
as a topological direct sum.

For each line define the compact slice
\[
 K_i=\{t\in\R:tv+b_i\in K\},
 \qquad J_i=\{g\in\Cin(\R):g|_{K_i}=0\}.
\]
Take the unique decomposition \(f=h+E(g_1,\ldots,g_q)\) from
\eqref{eq:ambient-split}.  Since \(h\) vanishes on every line and
\(E(g)|_{b_i+\R v}=g_i\), we have
\[
 f|_K=0 \quad\Longleftrightarrow\quad
 g_i|_{K_i}=0\ \text{ for every }i.
\]
Consequently,
\begin{equation}
 \I_K=\I_M\oplus E\!\left(\bigoplus_{i=1}^qJ_i\right)
 \label{eq:ideal-split}
\end{equation}
topologically.

It remains only to read \eqref{eq:ideal-split}.  By
Lemma~\ref{lem:permanence}, the complemented subspace \(\I_M\) of
\(\Cin(\R^d)\) has \(\Om\).  Each \(J_i\) is a closed ideal of
\(\Cin(\R)\), so it has
\(\Om\) by the source theorem \cite[Theorem 2.1]{FJW}.  Property \(\Om\)
is invariant under topological isomorphism.  The finite-sum part of
Lemma~\ref{lem:permanence} therefore shows that both summands in
\eqref{eq:ideal-split}, and hence \(\I_K\), have \(\Om\). \qed

\section*{Why this is only partial}

The proof works whenever the compact set lies in a carrier whose restriction
map has a continuous linear section and whose intrinsic vanishing ideal is
known to have \(\Om\).  Arbitrary compact subsets of \(\R^d\) need not admit
such a one-dimensional carrier.  A direct neighborhood cutoff does not close
the gap: zero values on \(K\) do not control normal derivatives, so cutting a
function off near \(K\) can fail even in \(C^1\).

The source's separate question about stability of the smooth extension
property under unions also remains open.  A naive fiber-product proof would
require a controlled decomposition of functions vanishing on \(K\cap L\)
into \(\I_K+\I_L\), and that algebraic equality can fail at accumulating
intersections.  The durable attempt note records eight distinct full-upgrade
routes and their obstructions.

\section*{Reproducible check and literature boundary}

The script \path{code/verify_parallel_line_splitting.py} constructs exact
transverse Lagrange functions for three lines, verifies \(RE=I\), and checks
the restriction identities on 303 samples.  This tests the explicit splitting,
not the abstract \(\Om\) stability facts.

No duplicate was found in the run indexes.  Exact-phrase searches and the
later citations of the published source were checked through 11 August 2026.
The only later citation returned by Semantic Scholar was a 2025 survey of
Whitney-jet spaces, which does not settle Conjecture 2.3.  This is a bounded
novelty check, not an exhaustive priority claim.

\begin{thebibliography}{9}
\bibitem{FJW}
L. Frerick, E. Jord\'a, and J. Wengenroth,
\emph{Extension operators for smooth functions on compact subsets of the
reals}, Math. Z. 295 (2020), 1537--1552; arXiv:1611.06808.
Included as \texttt{source\_paper.pdf}.

\bibitem{Vogt}
D. Vogt, \emph{Splitting of exact sequences of Fr\'echet spaces in the
absence of continuous norms}, J. Math. Anal. Appl. 297 (2004), 812--832.
Included under \texttt{references/} for the surrounding invariant and
splitting framework; the permanence facts used above are proved directly.
\end{thebibliography}

\end{document}
